\documentclass[aspectratio=169,14pt]{beamer}
\usetheme{default}
\usepackage{fontspec}
\setmainfont{texgyrepagella-regular.otf}[BoldFont=texgyrepagella-bold.otf,ItalicFont=texgyrepagella-italic.otf,BoldItalicFont=texgyrepagella-bolditalic.otf]
\setsansfont{texgyreheros-regular.otf}[BoldFont=texgyreheros-bold.otf,ItalicFont=texgyreheros-italic.otf,BoldItalicFont=texgyreheros-bold.otf,BoldItalicFeatures={FakeSlant=0.2}]
\setmonofont{lmmono10-regular.otf}[Scale=0.85]
\setbeamertemplate{navigation symbols}{}
\usepackage{pgfpages}
\newif\ifpresenter
\ifdefined\PresenterVersion\presentertrue\else\presenterfalse\fi
\ifpresenter
  \setbeameroption{show notes on second screen=right}
  \setbeamertemplate{note page}[plain]
\fi
\usepackage{amsmath,amssymb,bm,booktabs,array,graphicx,xcolor}
\usepackage{tikz,pgfplots,adjustbox,listings,xurl}
\pgfplotsset{compat=1.18}
\usetikzlibrary{arrows.meta,positioning,calc,fit,shapes.geometric,decorations.pathreplacing,patterns}
% UIUC official brand colors, verified 2026-10-23.
% https://brand.illinois.edu/visual-identity/color
\definecolor{IlliniOrange}{HTML}{FF5F05}
\definecolor{IlliniBlue}{HTML}{13294B}
\definecolor{Storm}{HTML}{707372}
\definecolor{LightStorm}{HTML}{C6C7C6}
\definecolor{Industrial}{HTML}{1D58A7}
\definecolor{Patina}{HTML}{007E8E}
\tikzset{>=Stealth, every node/.style={font=\small,text=IlliniBlue}, box/.style={draw=IlliniBlue,rounded corners=2pt,align=center,inner sep=6pt,execute at begin node=\hyphenpenalty10000\relax}, orangebox/.style={box,draw=IlliniOrange,very thick}, edge/.style={->,thick,IlliniBlue}, accent/.style={->,very thick,IlliniOrange}}
\pgfplotsset{every axis/.append style={width=13.5cm,height=6.5cm,axis line style={IlliniBlue},tick label style={font=\small,text=IlliniBlue},label style={font=\small,text=IlliniBlue},legend style={font=\small,draw=none,fill=white},grid style={LightStorm},major grid style={LightStorm},axis lines=left,clip=false}, every axis plot/.append style={thick}}

\definecolor{slidepaper}{RGB}{234,220,201}
\setbeamercolor{background canvas}{bg=slidepaper}
\setbeamercolor{normal text}{fg=IlliniBlue}
\setbeamercolor{frametitle}{fg=IlliniBlue}
\setbeamercolor{title}{fg=IlliniBlue}
\setbeamercolor{item}{fg=IlliniOrange}
\setbeamerfont{frametitle}{size=\large,series=\bfseries}
\setbeamerfont{title}{size=\Large,series=\bfseries}
\setbeamertemplate{itemize item}{\color{IlliniOrange}$\bullet$}
\setbeamertemplate{itemize subitem}{\color{IlliniOrange}\footnotesize$\circ$}
\setbeamercolor{orangerule}{bg=IlliniOrange}
\setbeamertemplate{headline}{\begin{beamercolorbox}[wd=\paperwidth,ht=0.1cm,dp=0.1cm]{orangerule}\end{beamercolorbox}}
\setbeamertemplate{footline}{\begin{beamercolorbox}[wd=\paperwidth,ht=0.1cm,dp=0.1cm]{orangerule}\end{beamercolorbox}}
\setbeamersize{text margin left=0.65cm,text margin right=0.65cm}
\setlength{\leftmargini}{1.1em}
\newenvironment{slidebullets}{\begin{itemize}\setlength\itemsep{0.8em}}{\end{itemize}}
\pgfplotsset{every axis/.append style={legend style={fill=slidepaper}}}
\lstdefinelanguage{CoursePython}{keywords={from,import,for,in,if,else,def,return,assert,True,False,None},sensitive=true,morecomment=[l]{\#},morestring=[b]',morestring=[b]"}
\lstset{language=CoursePython,basicstyle=\ttfamily\small,keywordstyle=\color{IlliniBlue}\bfseries,commentstyle=\color{Storm},stringstyle=\color{IlliniBlue},showstringspaces=false,breaklines=true,columns=fullflexible,keepspaces=true,frame=none,aboveskip=0.7em,belowskip=0.7em}
\newcommand{\E}{\mathbb E}
\newcommand{\Var}{\operatorname{Var}}
\newcommand{\Cov}{\operatorname{Cov}}
\hypersetup{colorlinks=true,urlcolor=IlliniBlue,linkcolor=IlliniBlue}

\title{Marketing Analytics}
\subtitle{Project Evidence and Midterm Presentations}
\author{By Nathan Yang}
\date{Session 08}
\begin{document}
\begin{frame}
\titlepage
\end{frame}
\begin{frame}{A Decision and Its Evidence}
\usebeamercolor[fg]{normal text}
\vfill
\begin{slidebullets}
\item State the action a manager can change.
\item Define the outcome, population, and horizon.
\item Identify which evidence could change the proposed action.
\end{slidebullets}
\vfill
\note{A presentation should make the decision assessable without requiring the audience to reconstruct the analyst's workflow. A precise unresolved question is useful when the current evidence does not yet support an action.}
\end{frame}
\begin{frame}{A Traceable Data Table}
\usebeamercolor[fg]{normal text}
\vfill
\begin{slidebullets}
\item Name the original unit and the analysis unit.
\item Reconcile rows and totals through joins and restrictions.
\item Show one manually checked record or aggregate.
\end{slidebullets}
\vfill
\note{A source-to-output audit distinguishes data processing choices from modeling choices. An unexplained change in the population can dominate a later coefficient or chart.}
\end{frame}
\begin{frame}{Evidence Categories}
\usebeamercolor[fg]{normal text}
\vfill
\begin{slidebullets}
\item Descriptive patterns summarize observed records.
\item Prediction evaluates outcomes on held-out observations.
\item Causal estimates require a justified counterfactual comparison.
\item Scenarios calculate consequences under imposed inputs.
\end{slidebullets}
\vfill
\note{These categories can appear in the same project. Each claim should retain the qualifications appropriate to its evidence.}
\end{frame}
\begin{frame}{A Worked Presentation Case}
\usebeamercolor[fg]{normal text}
\vfill
\begin{slidebullets}
\item A paid-plan scenario has high predicted uptake.
\item The underlying data contain food logging and no observed payments.
\item The decision requires additional evidence on price, retention, and margin.
\end{slidebullets}
\vfill
\note{Observed logging can motivate a segmentation or interface decision. Paid uptake from an imposed logit model is a scenario. A stated-choice study or pilot could supply some of the missing evidence.

M4, Section 4.11, Marketing Analytics Lecture Notes. \url{../06_Lecture_Notes/modules/M04_Conjoint_Segmentation_and_CLV.md}}
\end{frame}
\begin{frame}{An Interpretable Figure}
\usebeamercolor[fg]{normal text}
\vfill
\begin{slidebullets}
\item Use a comparison that answers the stated question.
\item Label units, denominator, and uncertainty.
\item State whether plotted quantities are observed, estimated, or imposed.
\end{slidebullets}
\vfill
\note{A figure should be independently readable. A chart of customer means and an interval for the overall mean require separate explanations of distribution and uncertainty.}
\end{frame}
\begin{frame}{A Reproducible Calculation}
\usebeamercolor[fg]{normal text}
\vfill
\begin{slidebullets}
\item Run from a clean session using documented inputs.
\item Show a counterexample that checks a consequential choice.
\item Keep complete code and results linked to the claim.
\end{slidebullets}
\vfill
\note{The most informative check is one that would fail under a plausible analytical error. Re-running identical code is useful but does not independently establish a valid interpretation.}
\end{frame}
\begin{frame}{Presentation Discussion}
\usebeamercolor[fg]{normal text}
\vfill
\begin{slidebullets}
\item What is the strongest supported claim?
\item Which assumption most affects the decision?
\item What evidence should be collected next?
\end{slidebullets}
\vfill
\note{Feedback should identify a concrete comparison or data requirement. A more complicated model is useful only if it addresses a stated limitation in the decision evidence.}
\end{frame}
\begin{frame}{A Revision Decision}
\usebeamercolor[fg]{normal text}
\vfill
\begin{slidebullets}
\item Prioritize the error with the largest consequence for the action.
\item Specify the additional data or comparison needed.
\item Preserve a reproducible baseline for the revised analysis.
\end{slidebullets}
\vfill
\note{Examples include fixing a multiplied join, changing a leakage-prone split, or measuring actual contribution. These changes can matter more than polishing an already clear chart.}
\end{frame}
\begin{frame}{A Customer-Outcome Audit}
\usebeamercolor[fg]{normal text}
\vfill
\begin{slidebullets}
\item Identify who benefits and who can bear a cost.
\item State whether the optimized metric can conflict with customer value.
\item Name the relevant jurisdiction and verify the rule's scope.
\end{slidebullets}
\vfill
\note{These questions organize an evidence discussion. They do not establish a graded assignment or replace the eventual project brief.

ETHICS\_AND\_REGULATION.md.}
\end{frame}
\begin{frame}{Sensitivity and Next Steps}
\usebeamercolor[fg]{normal text}
\vfill
\begin{slidebullets}
\item Report how the action changes across plausible inputs.
\item Identify a decision threshold rather than only a point estimate.
\item Select a feasible test that reduces consequential uncertainty.
\end{slidebullets}
\vfill
\note{A recommendation can be conditional. Explicitly stating the value or response threshold makes the condition useful for a manager rather than an abstract caveat.}
\end{frame}
\end{document}
